\chapter{Matching Algorithms and Implied Spreads}\label{ch:m1:matching-algorithms-and-implied-spreads}

Two traders bid the same price for the same future. A sell order arrives,
smaller than their combined size. In the equity-index future the trader who
bid first is filled completely before the other receives a single lot. In
the short-term interest-rate future next to it, on the same exchange, the two
share the sale in proportion to the sizes they showed, and the one who
arrived first has no advantage at all. The first rule pays for speed; the
second pays for size, and invites everyone to show more than they want. The
allocation rule is a few lines in the exchange's documentation and it decides
what kind of firm can make markets in the product. A second mechanism,
invisible on a screen, links the books of different expiries: the best bid
in a \omterm{def:m1:matching-algorithms-and-implied-spreads:calendar}{calendar spread} is often an order nobody entered.

\section{Price-time priority}

\begin{definition}[Price-time priority]\label{def:m1:matching-algorithms-and-implied-spreads:fifo}
Under \emph{price-time priority}\index{price-time priority} (first in, first
out, FIFO) an incoming order trades against resting orders at the best
price in the order of their priority timestamps, oldest first, each filled
completely before the next receives anything.
\end{definition}

An order's timestamp is set when it enters the book. On most exchanges it is
renewed, and the place in the queue lost, when the order's price is changed
or its quantity increased; reducing the quantity keeps the place. Under FIFO
a resting order has two attributes that matter: its price and its position
in the queue. At a large tick (\cref{fig:m1:futures-contracts-and-exchanges:tickbp})
the price rarely changes, queues are long, and position is the whole game:
being near the front means being filled by small, uninformed orders; being
at the back means being filled only when a large order sweeps the level,
which is when the price is about to move against the fill. The economics of
the queue are developed in One Quant Book 10.

\section{Pro rata and its variants}

\begin{definition}[Pro-rata allocation]\label{def:m1:matching-algorithms-and-implied-spreads:prorata}
Under \emph{pro-rata allocation}\index{pro-rata allocation} an incoming order
is shared among all resting orders at the best price in proportion to their
sizes. Allocations are rounded down to whole lots, allocations below a
product-specific minimum may be dropped, and what rounding leaves over is
distributed by a secondary rule, often time priority.
\end{definition}

\begin{definition}[Top-order allocation and lead market maker]\label{def:m1:matching-algorithms-and-implied-spreads:top}
A \emph{top-order allocation}\index{top-order allocation} gives the order
that first improved the market, establishing a new best price, a priority
allocation before any other rule applies. A \emph{lead market
maker}\index{lead market maker} is a firm which, in exchange for quoting
obligations, receives a set percentage of each incoming order at the best
price before the rest is allocated.
\end{definition}

Real algorithms are sequences of such steps. One family in use: the top
order first, up to a percentage of the aggressor; then \omterm{def:m1:matching-algorithms-and-implied-spreads:top}{lead market makers};
then a percentage of what remains by time; then the rest pro rata with a
minimum allocation; then leftovers by time.

\begin{dated}{2026-09}{Which products use which rule}\label{dat:m1:matching-algorithms-and-implied-spreads:which}
\textbf{CME Group} documents nine algorithms, among them FIFO, FIFO with
\omterm{def:m1:matching-algorithms-and-implied-spreads:top}{lead market maker}, Pro Rata, Allocation (top order, then pro rata with a
minimum of two lots, then FIFO), Threshold Pro Rata (with and without a \omterm{def:m1:matching-algorithms-and-implied-spreads:top}{lead
market maker}) and Configurable, a split of FIFO and pro rata with set
percentages. A data vendor's survey of the exchange's definitions assigns
FIFO to the equity-index, Treasury-note and crude-oil outrights, about 70\%
of volume; Configurable to two-year notes, federal funds and the grains,
about 13\%; Allocation to three-month SOFR outrights, about 10\%; and
Threshold Pro Rata with a \omterm{def:m1:matching-algorithms-and-implied-spreads:top}{lead market maker} to options on the ten-year note.
\textbf{Eurex} supports time, pro-rata and time-pro-rata allocation, the last
giving older orders more than their proportional share; on 30 March 2026 it
moved its money-market futures, the three-month Euribor among them, from time
to time-pro-rata allocation.
\end{dated}

The pattern is economic. Products whose price moves many ticks a day keep
FIFO: the tick is small, improving the price is the way to get priority, and
queues are short. Short-term interest-rate futures move a tick or two a day:
under FIFO the queue at each price would be days long and nobody but the
first arrival would bother to quote. Sharing fills pro rata keeps many firms
quoting size at a price that does not move.

\begin{omfigure}
\begin{tikzpicture}
\begin{axis}[width=11cm,height=5.2cm,ybar,bar width=9pt,
  ylabel={lots received},
  symbolic x coords={A,B,C,D},xtick=data,
  xticklabels={{A: 10, first}, {B: 200}, {C: 40, LMM}, {D: 250, last}},
  tick label style={font=\small},label style={font=\small},xticklabel style={font=\footnotesize\strut},
  ymin=0,ymax=105,ymajorgrids,grid style={gray!25},enlarge x limits=0.15,
  nodes near coords,nodes near coords style={font=\scriptsize},
  legend columns=3,legend style={font=\footnotesize,at={(0.5,-0.2)},anchor=north,draw=none,fill=none,/tikz/every even column/.append style={column sep=8pt}},
  table/col sep=comma]
\addplot[fill=omDef!70,draw=omDef] table[x=order,y=fifo]{figdata/markets-1/19-matching-algorithms-and-implied-spreads/alloc.csv};
\addplot[fill=omThm!60,draw=omThm] table[x=order,y=pro_rata]{figdata/markets-1/19-matching-algorithms-and-implied-spreads/alloc.csv};
\addplot[fill=omExo!50,draw=omExo] table[x=order,y=split]{figdata/markets-1/19-matching-algorithms-and-implied-spreads/alloc.csv};
\legend{FIFO,pro rata (minimum 2),split: top 10\% / LMM 10\% / 40\% by time / rest pro rata}
\end{axis}
\end{tikzpicture}

\omcaption{One hundred lots arrive at a level holding four orders, listed in
time priority with their sizes. The same aggressor, the same book, three
rules, three different sets of winners. Data: the chapter's
build.}\label{fig:m1:matching-algorithms-and-implied-spreads:alloc}
\end{omfigure}

\begin{proposition}[Pro rata rewards shown size]\label{prop:m1:matching-algorithms-and-implied-spreads:oversize}
Under pro rata without minimum, an order of size $s$ in a level whose other
orders total $Q$ receives $\lfloor as/(Q+s)\rfloor$ of an aggressor of size $a
\le Q+s$, increasing in $s$ and independent of arrival time. A trader who
wants $w$ lots from a typical aggressor $a$ must show about $s = wQ/(a-w)$,
and is exposed to receiving up to $s$ when a larger aggressor arrives.
\end{proposition}

\begin{example}[Eighty-seven to get eight]\label{ex:m1:matching-algorithms-and-implied-spreads:87}
A level holds 1\,000 lots and typical aggressors are 100 lots. To receive 8
the trader shows $8 \times 1\,000/92 = 87$ lots (86 would receive 7). If a
1\,000-lot order arrives instead, the trader receives 80: ten times the
intention,
at the moment the price is most likely to move through the level.
Displayed size in a pro-rata book therefore overstates the quantity anyone
wishes to trade, and much of it disappears when a large order appears.
\end{example}

\begin{omfigure}
\begin{tikzpicture}
\begin{axis}[width=10cm,height=4.6cm,
  xlabel={size shown by the newcomer (lots)},ylabel={lots received},
  tick label style={font=\small},label style={font=\small},
  xmin=0,xmax=400,ymin=0,ymax=32,grid=major,grid style={gray!25},table/col sep=comma]
\addplot[omDef,very thick,const plot] table[x=shown,y=filled]{figdata/markets-1/19-matching-algorithms-and-implied-spreads/oversize.csv};
\end{axis}
\end{tikzpicture}

\omcaption{Pro rata with a minimum allocation of two lots: what a newcomer
at the back of a 1\,000-lot level receives from a 100-lot aggressor, against
the size it shows. Below 21 lots it receives nothing. Data:
\cref{prop:m1:matching-algorithms-and-implied-spreads:oversize}.}\label{fig:m1:matching-algorithms-and-implied-spreads:oversize}
\end{omfigure}

\begin{omfigure}
\begin{tikzpicture}
\begin{axis}[width=10.5cm,height=4.8cm,
  xlabel={lots ahead in the queue},ylabel={expected fill (lots)},
  tick label style={font=\small},label style={font=\small},
  xmin=0,xmax=490,ymin=0,ymax=10.5,grid=major,grid style={gray!25},
  legend columns=2,legend style={font=\footnotesize,at={(0.5,-0.32)},anchor=north,draw=none,fill=none,/tikz/every even column/.append style={column sep=8pt}},
  table/col sep=comma]
\addplot[omDef,very thick] table[x=ahead,y=fifo]{figdata/markets-1/19-matching-algorithms-and-implied-spreads/position.csv};
\addplot[omThm,very thick,dashed] table[x=ahead,y=pro_rata]{figdata/markets-1/19-matching-algorithms-and-implied-spreads/position.csv};
\legend{FIFO,pro rata}
\end{axis}
\end{tikzpicture}

\omcaption{Expected fill of a 10-lot order in a 500-lot level when one
aggressor of random size (exponential, mean 120) arrives. Under FIFO the
front of the queue is worth four times the pro-rata share and the back
almost nothing; the curves cross with 170 lots ahead. Data: the tutorial's
simulation.}\label{fig:m1:matching-algorithms-and-implied-spreads:position}
\end{omfigure}

\section{Calendar spreads and implied prices}

\begin{definition}[Calendar spread]\label{def:m1:matching-algorithms-and-implied-spreads:calendar}
A \emph{calendar spread}\index{calendar spread} is an exchange-listed
instrument whose purchase is the simultaneous purchase of one expiry and
sale of another of the same product, at a quoted price difference. It has
its own order book, its own tick, and trades as one order with no risk of
being filled on one leg only.
\end{definition}

The convention used here: the spread is front minus back, so buying the
spread buys the front month and sells the back. Rolling a long position
forward is selling the spread.

\begin{definition}[Implied-in and implied-out]\label{def:m1:matching-algorithms-and-implied-spreads:implied}
An \emph{implied-in}\index{implied-in} order is a spread order which the
matching engine derives from orders resting in the outright books: a bid in
the front month and an offer in the back month imply a bid in the spread at
the difference of their prices, for the smaller of their sizes. An
\emph{implied-out}\index{implied-out} order is an outright order derived from
a spread order and an outright order in the other leg. Implied orders are
real: they can be hit, and the engine then executes all the legs at once.
Implied orders do not trade against implied orders.
\end{definition}

\begin{proposition}[Implied prices]\label{prop:m1:matching-algorithms-and-implied-spreads:implied}
With $b_F, a_F$ and $b_B, a_B$ the best bids and offers in the front and back
months and $b_S, a_S$ those entered directly in the spread:
\begin{align*}
\text{implied-in:}\quad & b_S^{\text{imp}} = b_F - a_B, & a_S^{\text{imp}} &= a_F - b_B,\\
\text{implied-out, front:}\quad & b_F^{\text{imp}} = b_S + b_B, & a_F^{\text{imp}} &= a_S + a_B,\\
\text{implied-out, back:}\quad & b_B^{\text{imp}} = b_F - a_S, & a_B^{\text{imp}} &= a_F - b_S.
\end{align*}
The size of each implied order is the smaller of its components' sizes.
\end{proposition}
\begin{proof}
A seller of the spread sells the front and buys the back. The engine can
fill that seller with the front-month bidder and the back-month seller, who
together pay $b_F$ and charge $a_B$. The other five are the same argument applied
to each leg.
\end{proof}

\begin{omfigure}
\begin{tikzpicture}[font=\small,
  bk/.style={draw=omDef,rounded corners=2pt,minimum width=3.3cm,minimum height=1.5cm,align=center,fill=omDef!4},
  flow/.style={-{Stealth[length=2mm]},thick}]
\node[bk] (f) at (0,1.2) {\textbf{Front month}\\bid 10\,050 $\times$ 30\\ask 10\,052 $\times$ 25};
\node[bk] (b) at (0,-1.2) {\textbf{Back month}\\bid 10\,020 $\times$ 40\\ask 10\,023 $\times$ 10};
\node[bk,draw=omThm,fill=omThm!5,minimum width=5.4cm,minimum height=2.3cm] (s) at (8.0,0) {\textbf{Spread (front $-$ back)}\\direct: bid 28 $\times$ 5, ask 32 $\times$ 7\\implied: bid 27 $\times$ 10, ask 32 $\times$ 25\\[2pt]\emph{shown}: bid 28 $\times$ 5, ask 32 $\times$ 32};
\draw[flow,omThm] (f.east) -- node[above,sloped,font=\footnotesize]{bid $-$} (s.west|-0,0.35);
\draw[flow,omThm] (b.east) -- node[below,sloped,font=\footnotesize]{$-$ ask} (s.west|-0,-0.35);
\end{tikzpicture}

\omcaption{\omterm{def:m1:matching-algorithms-and-implied-spreads:implied}{Implied-in} orders. The spread's implied bid is the front bid
minus the back ask, $10\,050 - 10\,023 = 27$, for $\min(30, 10) = 10$ lots; its
implied ask is $10\,052 - 10\,020 = 32$ for 25. The book a trader sees is the
better of direct and implied on each side, sizes adding at equal
prices.}\label{fig:m1:matching-algorithms-and-implied-spreads:implied}
\end{omfigure}

Implied matching joins the liquidity of all expiries: a hedger rolling a
position meets, through the spread, everyone quoting either month. It also
means that an order entered in one book can trade because of an event in
another, that an order's effective queue position depends on the engine's
rule for ranking implied against direct orders at the same price, and that a
simulator which ignores implieds understates the liquidity of the spread and
misstates fills in the outrights.

\section{What the feed shows}

\begin{definition}[Market by price and market by order]\label{def:m1:matching-algorithms-and-implied-spreads:mbo}
A \emph{market-by-price}\index{market-by-price} feed publishes, for each
price level, the total quantity and the number of orders. A
\emph{market-by-order}\index{market-by-order} feed publishes every order
individually, with an identifier and its priority, and every change to it.
\end{definition}

With market by order a firm knows its own position in a FIFO queue exactly
and sees who is ahead in size; with market by price it must estimate its
position from the sequence of updates, guessing whether each decrease in
quantity was a cancellation (ahead or behind?) or a trade. Implied
quantities are usually published separately from the direct book, or not at
all beyond the best level. Reading these feeds is the subject of
\cref{ch:m1:reading-market-data}; what matters here is that the allocation
rule determines which feed is worth paying for.

\section{Tutorial: three allocation rules}\label{tut:m1:matching-algorithms-and-implied-spreads:alloc}

\begin{tutorial}
\textbf{Goal.} Implement FIFO, pro rata with a minimum, and a configurable
split; compare their fills; compute implied spread prices. \textbf{End
state:} the three data figures of this chapter.
\begin{tutsteps}
\item \textbf{FIFO}, and a helper that never over-fills an order.
\omcode{code/firm/match/firm_match.py}{18}{30}{Time priority: walk the queue.}\label{lst:m1:matching-algorithms-and-implied-spreads:fifo}
\item \textbf{Pro rata.} Integer arithmetic only: floors, a minimum, and
leftovers by time.
\omcode{code/firm/match/firm_match.py}{38}{51}{Pro-rata allocation on the open quantities.}\label{lst:m1:matching-algorithms-and-implied-spreads:prorata}
\item \textbf{\omterm{def:m1:matching-algorithms-and-implied-spreads:implied}{Implied-in}.}
\omcode{code/firm/match/firm_match.py}{88}{93}{The spread quote implied by two outright quotes.}\label{lst:m1:matching-algorithms-and-implied-spreads:implied}
\item \textbf{Property tests.} On three hundred random books every algorithm
must allocate exactly $\min(\text{aggressor}, \text{resting})$ lots and never
more than an order's size.
\end{tutsteps}
\textbf{What to change next.} Implement time-pro-rata: weight each order by
size times a decreasing function of its rank in time, and find the weights
for which the first order of \cref{fig:m1:matching-algorithms-and-implied-spreads:alloc}
receives the same as under the split rule.
\end{tutorial}

\section{Build: allocation and implieds}\label{bld:m1:matching-algorithms-and-implied-spreads:match}

\begin{build}
\bfield{Purpose} The matching engine of the miniature firm's simulated
exchange (One Quant Book 10) delegates one decision to this module: given a
level and an aggressor, who gets what. The backtester uses the same code, so
that simulated fills obey the rule of the product simulated.
\bfield{Interface} \texttt{Resting(oid, qty, lmm, top)} in time priority;
\texttt{fifo(book, qty)}, \texttt{pro\_rata(book, qty, min\_alloc)},
\texttt{configurable(book, qty, top\_pct, lmm\_pct, fifo\_pct, min\_alloc)},
each returning fills by order; \texttt{Quote}; \texttt{implied\_in(front,
back)}, \texttt{implied\_out\_front(spread, back)}, \texttt{best\_of(direct,
implied)}.
\bfield{Rules} Integers only. Conservation: fills sum to $\min(\text{qty},
\text{resting})$. No order is over-filled. Implied quantity never matches
implied quantity: \texttt{best\_of} combines a direct and an implied quote
and is never applied to two implied ones.
\bfield{Acceptance tests} \texttt{code/firm/match/tests/}: the chapter's
worked allocations; the split rule; conservation on random books; implied
prices with an empty side.
\bfield{Stretch} A butterfly (three expiries) and the second generation of
implieds it creates; decide, as an exchange must, where to stop.
\end{build}

\begin{omsources}
\item CME Group Client Systems Wiki, \emph{Supported Matching Algorithms},
\emph{CME Globex Matching Algorithms} and \emph{Implied Orders}.
\item Databento, ``CME matching algorithms explained'' (survey of
algorithms by product and share of volume).
\item Eurex, \emph{Matching principles}, and the circular on money-market
derivatives (allocation scheme from 30 March 2026).
\end{omsources}

\section{Exercises}

\begin{exercise}[$\star$]\label{exo:m1:matching-algorithms-and-implied-spreads:1}
The level of \cref{fig:m1:matching-algorithms-and-implied-spreads:alloc} (A 10,
B 200, C 40, D 250, in time order) receives an aggressor of 230 lots under
FIFO. Give the fills.
\end{exercise}

\begin{exercise}[$\star$]\label{exo:m1:matching-algorithms-and-implied-spreads:2}
Same level, aggressor of 60 lots, pro rata without minimum and leftovers by
time. Give the exact shares, the floors, the leftover and the fills.
\end{exercise}

\begin{exercise}[$\star$]\label{exo:m1:matching-algorithms-and-implied-spreads:3}
Front month 10\,050 bid for 30, offered at 10\,052 for 25; back month 10\,020
bid for 40, offered at 10\,023 for 10. Give the \omterm{def:m1:matching-algorithms-and-implied-spreads:implied}{implied-in} bid and offer of
the spread with their sizes.
\end{exercise}

\begin{exercise}[$\star\star$]\label{exo:m1:matching-algorithms-and-implied-spreads:4}
Under FIFO, which of these actions on a resting order lose its place in the
queue on most exchanges: reducing its size; increasing its size; changing
its price; cancelling and re-entering it? How would you design an order
manager that needs to \emph{increase} its size at a level without giving up
the priority it has?
\end{exercise}

\begin{exercise}[$\star\star$]\label{exo:m1:matching-algorithms-and-implied-spreads:5}
Verify \cref{ex:m1:matching-algorithms-and-implied-spreads:87}. Then suppose
every participant reasons the same way and multiplies its size by ten. What
happens to each one's fills, to the displayed depth, and to the information
in the displayed depth?
\end{exercise}

\begin{exercise}[$\star\star$]\label{exo:m1:matching-algorithms-and-implied-spreads:6}
A \omterm{def:m1:matching-algorithms-and-implied-spreads:mbo}{market-by-price} feed shows the best bid going from 480 lots in 12 orders to
465 lots in 11 orders, with no trade message. Your own 10-lot order is
somewhere in that level under FIFO. What happened, and what can and cannot
be said about your queue position? What would market by order add?
\end{exercise}

\begin{exercise}[$\star\star\star$]\label{exo:m1:matching-algorithms-and-implied-spreads:7}
\emph{Coding.} From \texttt{position.csv}: report the expected fill at the
front of the queue under FIFO, the pro-rata expected fill, their ratio, and
the number of lots ahead beyond which pro rata would have been better.
\end{exercise}

\begin{exercise}[$\star\star\star$]\label{exo:m1:matching-algorithms-and-implied-spreads:8}
\emph{Find the flaw.} ``The product is pro rata with a two-lot minimum. I
will quote 5 lots on each side at the best price, in levels of about 2\,000
lots, and collect my fair share of the flow, about 0.25\% of each trade.''
Typical aggressors are 100 lots. What will this trader collect, and why?
\end{exercise}

\section{Problem: The Implied Book}

\begin{problem}[{Weekend problem --- two expiries, one spread, one engine}]\label{pb:m1:matching-algorithms-and-implied-spreads:1}
Prices are in ticks. Front month (Z): bid 10\,050 $\times$ 30, ask 10\,052
$\times$ 25. Back month (H): bid 10\,020 $\times$ 40, ask 10\,023 $\times$ 10.
Orders entered directly in the Z--H spread: bid 28 $\times$ 5, ask 32 $\times$
7. The spread is front minus back.

\medskip\noindent\textbf{Part I --- \omterm{def:m1:matching-algorithms-and-implied-spreads:implied}{Implied-in}.}
\begin{enumerate}
\item Give the \omterm{def:m1:matching-algorithms-and-implied-spreads:implied}{implied-in} bid of the spread, with its size.
\item Give the \omterm{def:m1:matching-algorithms-and-implied-spreads:implied}{implied-in} ask, with its size.
\item Give the spread book a trader sees at the best level.
\item Why can the implied bid (27) not be above the direct ask (32) for
long? What would happen if it were?
\item Explain why implied orders must not match implied orders.
\end{enumerate}

\medskip\noindent\textbf{Part II --- A roll.}
A fund sells 12 spreads at market (it rolls a long position forward).
\begin{enumerate}[resume]
\item Which orders does it trade against, at which prices?
\item Give its average price.
\item Which trades print in the outright books, and at which prices?
\item Give the three best-level books after the trade.
\item The fund could have sold 12 Z and bought 12 H itself. Give the prices it
would have obtained from the displayed outright books, and the difference.
\end{enumerate}

\medskip\noindent\textbf{Part III --- \omterm{def:m1:matching-algorithms-and-implied-spreads:implied}{Implied-out} and a new order.}
Return to the initial books.
\begin{enumerate}[resume]
\item Give the \omterm{def:m1:matching-algorithms-and-implied-spreads:implied}{implied-out} bid and ask in the front month, from the direct
spread orders and the back month.
\item Do they improve the front month's displayed market?
\item A new offer of 20 lots at 10\,021 enters the back month. Give the new
\omterm{def:m1:matching-algorithms-and-implied-spreads:implied}{implied-in} bid of the spread.
\item Does the new order trade? What would it take?
\item A \omterm{def:m1:what-a-trading-firm-does:market-maker}{market maker} rests a bid of 50 at 28 in the spread. After the order
of question 13, where does he stand relative to the implied bid, and what
does that tell him?
\end{enumerate}

\medskip\noindent\textbf{Part IV --- Judgement.}
\begin{enumerate}[resume]
\item Your backtest of a front-month strategy ignores the spread book. Name
two ways its fills are wrong.
\item The outrights are FIFO. An \omterm{def:m1:matching-algorithms-and-implied-spreads:implied}{implied-out} order and a direct order rest at
the same price. Why does the engine's ranking rule between them matter to a
\omterm{def:m1:what-a-trading-firm-does:market-maker}{market maker}?
\item Why do exchanges publish implied quantity only at the best level or
two?
\item State the \emph{named result}: the implied best bid of the spread and
the displayed best bid, with sizes.
\item In one sentence: what does implied matching do for a hedger?
\end{enumerate}
\end{problem}

\section{Interview questions}

\begin{interviewq}[$\star$ \iqroles{trader, developer, researcher}]\label{iq:m1:matching-algorithms-and-implied-spreads:1}
Explain FIFO and pro-rata matching. Which products use which, and why?
\end{interviewq}

\begin{interviewq}[$\star$ \iqroles{developer, trader}]\label{iq:m1:matching-algorithms-and-implied-spreads:2}
What is an implied order?
\end{interviewq}

\begin{interviewq}[$\star\star$ \iqroles{trader, researcher}]\label{iq:m1:matching-algorithms-and-implied-spreads:3}
How does a \omterm{def:m1:what-a-trading-firm-does:market-maker}{market maker}'s behaviour differ between a FIFO and a pro-rata
product?
\end{interviewq}

\begin{interviewq}[$\star\star$ \iqroles{developer}]\label{iq:m1:matching-algorithms-and-implied-spreads:4}
Write \omterm{def:m1:matching-algorithms-and-implied-spreads:prorata}{pro-rata allocation} with a minimum lot size. What are the edge cases?
\end{interviewq}

\begin{interviewq}[$\star\star$ \iqroles{researcher, mle}]\label{iq:m1:matching-algorithms-and-implied-spreads:5}
You only have \omterm{def:m1:matching-algorithms-and-implied-spreads:mbo}{market-by-price} data. How do you estimate your queue position
in a FIFO book?
\end{interviewq}

\begin{interviewq}[$\star\star\star$ \iqroles{developer, researcher}]\label{iq:m1:matching-algorithms-and-implied-spreads:6}
You are building an exchange simulator for \omterm{def:m1:matching-algorithms-and-implied-spreads:calendar}{calendar spreads}. What does
implied matching force you to get right?
\end{interviewq}
